\documentclass[11pt,oneside]{memoir}
\usepackage{raeez-math-template}
\title{The nonadjacent collision and its source extension}
\author{}
\date{}
\begin{document}
\maketitle
\begin{abstract}
At nonzero Heisenberg level, the first-to-third collision cannot be represented by outputs supported on the adjacent diagonals.
Its minimal cyclic source module over regular position functions has length two in the normal direction.
Adjoining this module supplies the missing source differential on the open complement of the adjacent diagonals.
The resulting comparison takes values in the unextended chiral algebra.
An exact sequence computes the change in source cohomology globally.
It includes a length-four quotient on the small diagonal and an explicit failure of quasi-isomorphism.
An all-word variant uses the same added module and has a length-three small-diagonal quotient.
The new collision module has a specified right differential module presentation.
A separate binary calculation excludes a regular coefficient-compatible connection on the unchanged ordered symbols.
\end{abstract}
\chapter{The nonadjacent collision and its source extension}
\label{chap:bdnz-source}

\section{The native double-pole output}

Fix $k\in\mathbb C^\times$ and $X=\mathbb A^1_{\mathbb C}$.
The level remains this nonzero scalar throughout.
The Heisenberg states are
\[
 V_k=\mathbb C[x_1,x_2,\ldots],\qquad
 \mathbf1=1,\quad J=x_1,\quad
 T=\sum_{n\geq1}n x_{n+1}\partial_{x_n}.
\]
All states are even, have cohomological degree zero, and have zero differential.
Every state is a polynomial in finitely many variables.
The modes are $J_{-n}=x_n\cdot$, $J_n=kn\partial_{x_n}$ for $n\geq1$, and $J_0=0$.
The fields of monomials are the normally ordered products of the divided current derivatives.
In particular,
\begin{equation}\label{eq:bdnz-field}
 Y(J,t)J=\frac{k}{t^2}\mathbf1+(e^{tT}J)J,
 \qquad Y(\mathbf1,t)a=a.
\end{equation}
Normal ordering places multiplication operators before derivative operators.
Commuting these operators gives the finite Wick pairing formula and the usual Heisenberg vertex operations.

The right differential module is $\mathcal H_k=\omega_X\otimes V_k$, with
\[
 (fa\,dz)\partial_z=-(f'a+fTa)\,dz.
\]
Its chiral operation on a pole-valued input is
\[
 \mu(f(z,w)(a\,dz\boxtimes b\,dw))
   =\operatorname{pp}_{t}\bigl(f(y+t,y)Y(a,t)b\bigr),
 \qquad t=z-w,\quad y=w.
\]
The symbol $\operatorname{pp}_t$ denotes the finite principal part.
Multiplication occurs before passage to the quotient.
The normalization is
\begin{equation}\label{eq:bdnz-normal}
 \operatorname{pp}_t\frac{a(y)}{t^{r+1}}
 =\frac{a(y)\,dy}{r!}\delta_\Delta\partial_t^r,
 \qquad dz\wedge dw=dt\wedge dy.
\end{equation}
All normal derivatives are retained, and all sums of normal derivatives are finite.

Let $R=\mathbb C[z_1,z_2,z_3]$.
For a partition $\pi$ of $\{1,2,3\}$, let $\Delta_\pi:X^\pi\hookrightarrow X^3$ identify its blocks.
Let $j_\pi$ include the locus of distinct block positions.
The native three-point Chevalley--Cousin complex is
\begin{equation}\label{eq:bdnz-native}
 E_k=
 \bigoplus_{\pi}
 (\Delta_\pi)_!(j_\pi)_*j_\pi^*
 \left(\boxtimes_{B\in\pi}\mathcal H_k[1]\right).
\end{equation}
The closed direct image uses the right transfer module, including every normal derivative.
There is no additional geometric degree shift.
This is the carrier of \cite{BDnonadjacent}*{\S3.4.11, equation (3.4.11.1)}.
The separated summand has degree $-3$, the three pair summands have degree $-2$, and the small-diagonal summand has degree $-1$.

Write $[a|b|c]$ for the suspended states, each of degree $-1$.
The collision differential satisfies
\begin{equation}\label{eq:bdnz-jacobi-signs}
 Q[a|b|c]=[\mu_{12}(a,b)|c]-[\mu_{13}(a,c)|b]
                      +[\mu_{23}(b,c)|a],
 \qquad Q[a|b]=[\mu(a,b)].
\end{equation}
Reordering two suspended even states changes the sign.
Each map includes its coordinate permutation, transverse localization, and transfer module.
Chiral Jacobi gives $Q^2=0$.
There is no internal state differential in $E_k$.

Put
\[
 t=z_1-z_3,\qquad s=z_2-z_3,\qquad y=z_3,
 \qquad p=(z_1-z_2)(z_2-z_3)=s(t-s).
\]
The common orientation is $dz_1\wedge dz_2\wedge dz_3=dt\wedge ds\wedge dy$.
Fix the block order $13|2$ and the normal coordinate $t$.
The section
\begin{equation}\label{eq:bdnz-nu}
 \nu=k[\mathbf1\delta_{13}\partial_t|J_2]\in E_k^{-2}
\end{equation}
is nonzero.
The factor $k$ is part of the prescribed current operation.
It is not changed or set to zero.

\section{The support obstruction}

Let
\[
 U=X^3\setminus(\Delta_{12}\cup\Delta_{23}),
 \qquad A=R[p^{-1}].
\]
This open set meets $\Delta_{13}$ and excludes the small diagonal.
In particular, $s$ and $t-s$ are invertible on $U$.
The restricted native complex has only two nonzero partition summands:
the separated summand in degree $-3$ and the $13$ summand in degree $-2$.
The latter differential is zero because its next collision is absent from $U$.

\begin{theorem}\label{thm:bdnz-support}
Suppose an ordered three-letter source has
\[
 be_{123}=-\omega_{12}m_{12|3}-\omega_{23}n_{1|23}.
\]
No degree-zero $R$-linear cochain comparison to $E_k$ can satisfy both of the following requirements:
the separated component of $e_{123}$ is $[J_1|J_2|J_3]$,
and the images of its two multiplication faces are supported on $\Delta_{12}$ and $\Delta_{23}$, respectively.
Allowing further corrections supported on these adjacent diagonals or on their small intersection does not remove the obstruction.
\end{theorem}
\begin{proof}
Restrict the proposed chain equation to $U$.
All stated correction terms and both source-face images vanish there.
The $13$ component of the separated insertion differential is
\[
 Q[J_1|J_2|J_3]=-\nu\qquad\hbox{on }U.
\]
The minus sign is the middle unshuffle sign in \eqref{eq:bdnz-jacobi-signs}.
Formula~\eqref{eq:bdnz-field} gives the coefficient $k\delta_{13}\partial_t$.
This coefficient is nonzero at every point of $\Delta_{13}\cap U$.
The current $J_2$ is nonzero and the normal derivatives form the free transfer basis.
Thus the two sides of the chain equation differ by the nonzero section $-\nu$.
\end{proof}

In the unextended complex, degree alone also forbids a partial-diagonal correction to $e_{123}$ itself.
Its degree is $-3$, whereas every proper-diagonal state has degree $-2$.
The missing term must therefore occur in the source differential if the separated insertion and the unextended target are retained.

\section{The ordered source and its lowest cohomology}

Let $\mathcal O_3$ be the six total orders of $\{1,2,3\}$.
For nonempty $B\subset\mathcal O_3$, let $R_B$ invert $z_i-z_j$ exactly when the relative order of $i,j$ varies in $B$.
Let
\[
 \Omega_B=\mathbb C[u_\sigma,du_\sigma\mid\sigma\in B]
 /(\textstyle\sum u_\sigma-1,\sum du_\sigma),
\]
with exterior odd generators $du_\sigma$.
Let $C$ be the differential graded algebra of face-compatible families in $R_B\otimes\Omega_B$.
Its differential $d_C$ fixes all positions and sends $u_\sigma$ to $du_\sigma$.
Define
\[
 q_{ij}=\sum_{i\prec_\sigma j}u_\sigma,
 \qquad \omega_{ij}=\frac{dq_{ij}}{z_i-z_j}.
\]
On a component with no inverse for the denominator, the numerator vanishes and that component is zero.

For each ordered triple $\sigma=ijk$, let
\[
 e_\sigma=[a_i|a_j|a_k],\quad
 m_\sigma=[b_{ij}|a_k],\quad
 n_\sigma=[a_i|b_{jk}].
\]
The states $a_i$ have degree zero and $b_{ij}$ have degree $-1$.
Suspension lowers degree by one.
All three displayed words therefore have degree $-3$.
The six triples give eighteen distinct generators.
Their direct sum of free $C$-modules is denoted $O$.
Its differential is
\begin{equation}\label{eq:bdnz-ordered}
 \begin{gathered}
 be_\sigma=-\omega_{ij}m_\sigma-\omega_{jk}n_\sigma,
 \qquad bm_\sigma=bn_\sigma=0,\\
 b(cv)=d_Cc\,v+(-1)^{|c|}c\,bv.
 \end{gathered}
\end{equation}
These are the suspended bar signs for
$a_i a_j=\omega_{ij}b_{ij}$, with every other nonunit product involving a $b$ equal to zero.
The closed forms $\omega_{ij}$ give $b^2=0$.
The complex is the direct sum of six three-generator subcomplexes $O_\sigma$.
It has no degree below $-3$.

\begin{lemma}\label{lem:bdnz-lowest}
For each $\sigma=ijk$, put $t_{ij}=z_i-z_j$ and $t_{jk}=z_j-z_k$.
Then
\begin{equation}\label{eq:bdnz-zsigma}
 z_\sigma=t_{ij}t_{jk}e_\sigma
           +t_{jk}q_{ij}m_\sigma+t_{ij}q_{jk}n_\sigma
\end{equation}
is a cycle, and
\[
 H^{-3}(O)=\bigoplus_{\sigma\in\mathcal O_3}
 \left(Rz_\sigma\oplus Rm_\sigma\oplus Rn_\sigma\right).
\]
\end{lemma}
\begin{proof}
A closed zero-form in $C$ is a constant polynomial on the full simplex, with coefficient in the common position localization.
Its vertex restrictions belong to $R$.
Injectivity of every coefficient localization then shows that all its components equal one element of $R$.
Thus $\ker(d_C:C^0\to C^1)=R$.

A degree-$-3$ cycle in $O_\sigma$ has the form $fe_\sigma+gm_\sigma+hn_\sigma$, with
\[
 f\in R,\qquad d_Cg=f\omega_{ij},\qquad d_Ch=f\omega_{jk}.
\]
Restrict the first equation to the actual edge $\{ijk,jik\}$, with coordinate $u=q_{ij}$.
Its coefficient ring is $R[t_{ij}^{-1}]$ and both endpoints of $g$ lie in $R$.
Integration gives $g(1)-g(0)=f/t_{ij}$.
Consequently $t_{ij}$ divides $f$ in $R$.
The edge $\{ijk,ikj\}$ proves that $t_{jk}$ divides $f$.
These two distinct linear polynomials are relatively prime, so $f=t_{ij}t_{jk}a$ for some $a\in R$.

The displayed $z_\sigma$ is closed by $d_Cq_{ij}=t_{ij}\omega_{ij}$ and its analogous identity for $jk$.
Subtracting $a z_\sigma$ leaves two closed coefficients, hence two elements of $R$.
There are no incoming boundaries in degree $-3$.
The three leading coefficients show independence, proving the assertion in every summand.
\end{proof}

\section{The minimal nonadjacent source module}

Let $N=R\nu$ be the cyclic $R$-submodule of the $13$ summand generated by \eqref{eq:bdnz-nu}.
Regard it as a complex concentrated in degree $-2$.
The unit field in \eqref{eq:bdnz-field} gives $Q\nu=0$ globally: its remaining coefficient is regular in the outer variable.
Thus $N$ is a subcomplex of $E_k$.
Regular function multiplication preserves this statement by $R$-linearity of $Q$.

\begin{proposition}\label{prop:bdnz-minimal}
The map $R\to N$, $f\mapsto f\nu$, has kernel $(t^2)$.
In particular,
\begin{equation}\label{eq:bdnz-two-jets}
 N\cong R/(t^2),\qquad
 t\nu=k[\mathbf1\delta_{13}|J_2]\ne0,\qquad t^2\nu=0.
\end{equation}
This is the minimal cyclic $R$-module of source outputs whose chosen generator maps to the native double-pole coefficient.
\end{proposition}
\begin{proof}
Write a polynomial uniquely as $f=f_0(y,s)+t f_1(y,s)+t^2g$.
The transfer relation gives
\[
 f\nu=k f_0[\mathbf1\delta_{13}\partial_t|J_2]
       +k f_1[\mathbf1\delta_{13}|J_2].
\]
The two normal coefficients are independent and $k\ne0$.
Thus the image vanishes exactly when $f_0=f_1=0$.
Any cyclic module with a specified generator mapping to $\nu$ surjects onto this image.
In particular, the relation $t=0$ on the generator is incompatible with the specified image.
\end{proof}

Define a degree-one $R$-linear map $\beta:O\to N$ by
\begin{equation}\label{eq:bdnz-beta}
 \beta(fe_{123})=-f(123)\nu\quad(f\in C^0),
 \qquad \beta=0\quad\hbox{on every other coefficient degree and word}.
\end{equation}
Evaluation $f(123)$ is at the order vertex $123$.
The map kills $bO$, since a coefficient differential has positive degree and a multiplication face is mixed.
Hence $\beta b=0$.
Set
\begin{equation}\label{eq:bdnz-extension}
 O^+=O\oplus N,
 \qquad b^+(x,n)=(bx,\beta x).
\end{equation}
The degree of $\beta$ makes this a degree-one differential, and $(b^+)^2=0$.
In particular, the new source operation is $b^+e_{123}=be_{123}-\nu$.
It has support on the actual nonadjacent diagonal and carries both normal jets in \eqref{eq:bdnz-two-jets}.
This construction is an extension of cochain complexes over $R$.
It makes no claim to be a bar coalgebra or a multiplication on the original ordered state algebra.

\section{The exact global change in cohomology}

The inclusion and projection give a short exact sequence of $R$-complexes
\begin{equation}\label{eq:bdnz-exact}
 0\longrightarrow N\longrightarrow O^+
       \xrightarrow{\pi}O\longrightarrow0.
\end{equation}
Its connecting homomorphism $H^{-3}(O)\to N$ is induced by $\beta$.
For the selected word, Lemma~\ref{lem:bdnz-lowest} gives
\begin{equation}\label{eq:bdnz-connecting}
 \beta[z_{123}]=-p\nu=-s(t-s)\nu,
 \qquad \beta[m_\sigma]=\beta[n_\sigma]=0.
\end{equation}
It vanishes on $z_\sigma$ for $\sigma\ne123$.

\begin{theorem}\label{thm:bdnz-cohomology}
The image of $H^{-3}(\pi)$ is obtained from the basis in Lemma~\ref{lem:bdnz-lowest} by replacing $Rz_{123}$ with $t^2Rz_{123}$.
The map is injective, and its cokernel is isomorphic to $R/(t^2)$.
In the next degree there is an exact sequence
\begin{equation}\label{eq:bdnz-hminus2}
 0\longrightarrow R/(t^2,s(t-s))
 \longrightarrow H^{-2}(O^+)
 \xrightarrow{H^{-2}(\pi)}H^{-2}(O)\longrightarrow0.
\end{equation}
For all other degrees, $H^n(\pi)$ is an isomorphism.
The projection $\pi$ is not a quasi-isomorphism.
\end{theorem}
\begin{proof}
Multiplication by $p=s(t-s)$ is injective on $R/(t^2)$.
Indeed, if $a=a_0+t a_1$ modulo $t^2$, then
\[
 pa=-s^2a_0+t(sa_0-s^2a_1)\pmod {t^2}.
\]
The first coefficient forces $a_0=0$, and the second then forces $a_1=0$.
Thus the kernel of \eqref{eq:bdnz-connecting} consists exactly of the asserted $t^2$ multiples in the selected component and the other unchanged components.
The connecting image is $pN$, and $N/pN=R/(t^2,p)$.
The long exact sequence of \eqref{eq:bdnz-exact} gives every displayed assertion.
Its only potentially affected degrees are $-3$ and $-2$, since $N$ occurs only in degree $-2$.
The nonzero connecting map excludes a quasi-isomorphism.
\end{proof}

The new degree-$-2$ submodule is supported on the small diagonal $t=s=0$.
It is free of rank four over $\mathbb C[y]$, with basis
\begin{equation}\label{eq:bdnz-lengthfour}
 1,\ s,\ t,\ st,
 \qquad t^2=0,\quad s^2=st,\quad s^3=0.
\end{equation}
For independence, view the quotient first as a free module with basis $1,t$ over $\mathbb C[y,s]$.
Its additional relations are $s^2=st$ and $s^2t=0$.
They reduce every monomial to the four displayed classes.
For independence, use lexicographic order $s>t$.
The polynomials $t^2$ and $s^2-st$ form a Gr\"obner basis: their only $S$-polynomial is a multiple of $t^2$.
Their leading monomials are $t^2$ and $s^2$, so the four displayed normal monomials are independent.
Thus the added nonadjacent operation creates a specific small-diagonal cohomology contribution.
The exact sequence does not assert a splitting of \eqref{eq:bdnz-hminus2}.

\section{The native comparison away from adjacent collisions}

Localize \eqref{eq:bdnz-exact} at $p$.
Write $O_A=O\otimes_R A$ and $N_A=N\otimes_R A$.
On the selected sector, the element
\begin{equation}\label{eq:bdnz-open-cycle}
 e'=e_{123}+\frac{q_{12}}{z_1-z_2}m_{12|3}
              +\frac{q_{23}}{z_2-z_3}n_{1|23}
       =p^{-1}z_{123}
\end{equation}
is a genuine cycle in $O_A$.
Both denominators are units in $A$.
Its new boundary component is $\beta e'=-\nu$.

\begin{theorem}\label{thm:bdnz-native-map}
There is a cochain map
\[
 F^+:O_A^+\longrightarrow E_k|_U
\]
whose separated projection on the selected word is
$fe_{123}\mapsto f(123)[J_1|J_2|J_3]$.
It sends the new source module $N_A$ to its identical native collision module.
It sends all mixed words, all positive-degree coefficients of the selected word, and all other ordered-word sectors to zero.
On $U$, the cohomology projection satisfies
\[
 \operatorname{coker}H^{-3}(\pi_A)=N_A,
 \qquad H^n(\pi_A)\text{ is an isomorphism for }n\ne-3.
\]
\end{theorem}
\begin{proof}
Let $B_U$ be the prescribed separated vertex map, zero on the other terms.
It kills $bO_A$.
By \eqref{eq:bdnz-jacobi-signs},
\[
 QB_U(x)=\beta x\quad\hbox{in the native }13\hbox{ summand}.
\]
Set $F^+(x,n)=B_U(x)+n$.
The module $N_A$ is closed in the target.
Therefore
\[
 QF^+(x,n)=\beta x=F^+(bx,\beta x)=F^+b^+(x,n).
\]
The specified degrees and supports give the component assertions.

Localization is exact and commutes with cohomology.
Since $p$ is a unit in $A$, the connecting homomorphism is surjective by \eqref{eq:bdnz-connecting}.
The small-diagonal quotient in \eqref{eq:bdnz-hminus2} vanishes after localization.
The cohomology assertions follow from Theorem~\ref{thm:bdnz-cohomology}.
\end{proof}

This is a comparison for the specified ordered insertion on its actual open coefficient domain.
It does not assign the pure separated current to every other ordered word.
Relabeling constructs the same comparison for each ordered word on the complement of its own two adjacent diagonals.
Compatibility among those comparisons on their intersections requires additional equations.

The class $[e']$ in the original source cannot lift to a cycle of the extended source.
Its connecting image is the nonzero element $-\nu$.
More generally, no quasi-isomorphic source replacement can preserve this projected insertion and give a lift to the native restricted complex.
Indeed, the connecting map for the two-term target sends its separated cohomology class to $-\nu$.
This composite is nonzero on $[e']$ and therefore nonzero in the derived category of $A$-modules.
A lift through the target would make that composite zero.
The extension \eqref{eq:bdnz-extension} changes the source cohomology exactly where this derived obstruction requires a change.

\section{All ordered insertions on the same open set}

The same cyclic module can receive the native $13$ boundary of every ordered word on $U$.
This gives a second extension which assigns pure separated currents to all six words.
Let $\varepsilon_\sigma$ be the sign of the permutation $\sigma=ijk$ of $123$.
Suspension gives
$[J_i|J_j|J_k]=\varepsilon_\sigma[J_1|J_2|J_3]$ in the fixed separated summand.
Define
\[
 \beta_{\mathrm{all}}(fe_\sigma)=-\varepsilon_\sigma f(\sigma)\nu,
 \qquad f\in C^0,
\]
and set it to zero on mixed words and positive coefficient degrees.
It again kills $bO$.
Put $O^{\mathrm{all}}=O\oplus N$ with differential
$b^{\mathrm{all}}(x,n)=(bx,\beta_{\mathrm{all}}x)$.
Only one new cyclic module is adjoined.

\begin{theorem}\label{thm:bdnz-all-words}
The formula
\[
 F^{\mathrm{all}}(fe_\sigma)=f(\sigma)[J_i|J_j|J_k],
 \qquad F^{\mathrm{all}}|_{N_A}=\mathrm{id},
\]
with zero images on all mixed words and positive-degree coefficients, defines a cochain map
$O^{\mathrm{all}}\otimes_R A\to E_k|_U$.
Globally, the connecting homomorphism sends $z_\sigma$ to $\lambda_\sigma\nu$, where
\[
 \begin{array}{c|rrrrrr}
 \sigma&123&132&213&231&312&321\\ \hline
 \lambda_\sigma\bmod t^2&s^2-st&-st&st&st&-st&st-s^2.
 \end{array}
\]
It kills every mixed basis class.
The induced cohomology sequence is
\begin{equation}\label{eq:bdnz-all-cohomology}
 \begin{gathered}
 0\longrightarrow H^{-3}(O^{\mathrm{all}})
 \longrightarrow H^{-3}(O)
 \xrightarrow{\beta_{\mathrm{all}}}(s^2,st)\,R/(t^2)
 \longrightarrow0,\\
 0\longrightarrow R/(t^2,s^2,st)
 \longrightarrow H^{-2}(O^{\mathrm{all}})
 \longrightarrow H^{-2}(O)\longrightarrow0.
 \end{gathered}
\end{equation}
All other cohomology groups are unchanged.
The small-diagonal quotient is free of rank three over $\mathbb C[y]$, with basis $1,s,t$.
\end{theorem}
\begin{proof}
On $U$, the differential of each ordered separated insertion is
$-\varepsilon_\sigma\nu$.
Its coefficient differential and both old multiplication faces have zero image.
The added source term is exactly $\beta_{\mathrm{all}}$, and the new module is closed in the target.
This proves the chain equation on every source summand.

Lemma~\ref{lem:bdnz-lowest} gives
$\lambda_\sigma=-\varepsilon_\sigma t_{ij}t_{jk}$.
Substitute $z_1=y+t$, $z_2=y+s$, and $z_3=y$, and reduce modulo $t^2$.
This gives the six displayed coefficients.
They generate precisely the ideal $(s^2,st)$ in $R/(t^2)$.
The long exact cohomology sequence gives \eqref{eq:bdnz-all-cohomology}.
The quotient has the monomial ideal $(t^2,s^2,st)$, whose standard monomials are $1,s,t$.
They prove the asserted rank and support.
\end{proof}

The kernel in the first sequence is explicit: a lowest source class
$\sum_\sigma(a_\sigma z_\sigma+b_\sigma m_\sigma+c_\sigma n_\sigma)$ lifts exactly when
$\sum_\sigma\lambda_\sigma a_\sigma\in(t^2)$.
On $U$, $p$ is a unit, so the connecting map is surjective onto $N_A$ and the small-diagonal quotient disappears.
The cohomology projection remains noninvertible in degree $-3$.
All old mixed images in this native comparison are zero.
Thus this extension has different mixed comparison data from a comparison to a differential graded state enlargement.
It does not identify either extended source with the original ordered source.

\section{The differential module of the new collision}

The cyclic module $N$ is minimal over regular position functions.
Its closure under right differential operators is larger and has a direct presentation.
Let $\mathcal D_R$ be the Weyl algebra in $t,s,y$, with $[\partial_t,t]=[\partial_s,s]=[\partial_y,y]=1$ and all other coordinate commutators zero.
Let $N_{\mathcal D}=\nu\mathcal D_R$ inside the native $13$ summand.

\begin{proposition}\label{prop:bdnz-weyl}
As a right differential module,
\begin{equation}\label{eq:bdnz-weyl-presentation}
 N_{\mathcal D}\cong
 \mathcal D_R\big/
 \left(t^2\mathcal D_R+(t\partial_t-1)\mathcal D_R
                 +(\partial_y-\partial_s)\mathcal D_R\right).
\end{equation}
It has independent finite normal forms with coefficients in $\mathbb C[y,s]$ and states
\[
 [\mathbf1\delta_{13}\partial_t^r|T^nJ_2],
 \qquad r,n\geq0.
\]
It is a closed subcomplex of the native target in degree $-2$.
\end{proposition}
\begin{proof}
The relation $\nu t^2=0$ was proved above.
Since $\nu t=k[\mathbf1\delta_{13}|J_2]$, right differentiation gives
$\nu t\partial_t=\nu$.
Translation of the vacuum is zero.
In these coordinates, both $\partial_s$ and $\partial_y$ act on the constant-coefficient generator as minus translation of $J_2$.
Hence $\nu(\partial_y-\partial_s)=0$.

Order the Weyl monomials with the tangent coordinate functions on the right.
The third relation eliminates $\partial_y$ from the generator.
The first two relations reduce the normal factors to $t$ and $\partial_t^r$ for $r\geq0$.
Thus the quotient is spanned by
\[
 \nu t\partial_s^n f(y,s),\qquad
 \nu\partial_t^r\partial_s^n f(y,s).
\]
Their images have normal orders zero and $r+1$, respectively, and state factors $(-T)^nJ$.
The states $T^nJ=n!x_{n+1}$ are independent.
The normal transfer basis is also independent, so no further relation exists.
This proves the presentation and normal forms.
Finally $Q\nu=0$ and right differential linearity give $QN_{\mathcal D}=0$.
\end{proof}

Induce the original source by $\operatorname{Ind}(O)=O\otimes_R\mathcal D_R$, with differential $b\otimes1$.
The degree-one map
\[
 \beta_{\mathcal D}(x\otimes P)=\beta(x)P
\]
has values in $N_{\mathcal D}$ and satisfies $\beta_{\mathcal D}(b\otimes1)=0$.
The extension by $N_{\mathcal D}$ therefore defines a complex of right differential modules.
On $U$, its comparison to $E_k|_U$ sends
\[
 (x\otimes P,n)\longmapsto B_U(x)P+n.
\]
Balancing over $R$ and right differential linearity prove the chain equation from Theorem~\ref{thm:bdnz-native-map}.
The relations in \eqref{eq:bdnz-weyl-presentation} are explicitly killed by this map.
They specify the differential closure of the added collision source.
They impose no unproved translation relation on the original ordered symbols.
Replacing $\beta$ and $B_U$ by the all-word maps gives the same right differential module construction for $O^{\mathrm{all}}$.

\section{A binary obstruction to a regular coefficient connection}

The original binary source has a stronger restriction than a missing choice of connection.
Let $R_2=\mathbb C[y,t]$, $L_2=R_2[t^{-1}]$, and
\[
 C_2^0=\{f\in L_2[u]:f(0),f(1)\in R_2\},
 \qquad C_2^1=L_2[u]du,\qquad d_Cf=f' du.
\]
Let $M=C_2e\oplus C_2h$, with $|e|=|h|=-2$, and
\[
 be=-\frac{du}{t}h,\qquad bh=0,
 \qquad b(cm)=d_Cc\,m+(-1)^{|c|}c\,bm.
\]
A coefficient-compatible normal connection means a degree-zero operator $\nabla_t$ satisfying
\[
 \nabla_t(cm)=(\partial_tc)m+c\nabla_tm,
 \qquad \partial_tu=0,\quad\partial_tdu=0.
\]
Regularity requires $\nabla_te,\nabla_th\in M^{-2}$ with their coefficients in $C_2^0$.

\begin{proposition}\label{prop:bdnz-no-connection}
No regular coefficient-compatible normal connection on $M$ commutes with $b$.
This statement does not depend on the Heisenberg level.
\end{proposition}
\begin{proof}
A degree-$-2$ cycle $fe+gh$ satisfies $\partial_uf=0$ and $\partial_ug=f/t$.
Endpoint regularity then gives
\[
 Z^{-2}(M)=R_2(te+uh)\oplus R_2h.
\]
Suppose the stated connection exists, and write $\nabla_te=Ae+Bh$ with $A,B\in C_2^0$.
The element $te+uh$ is closed, so its covariant derivative must also be closed.
Coefficient compatibility gives
\[
 \nabla_t(te+uh)=e+t\nabla_te+u\nabla_th.
\]
At $u=0$, its $e$ coefficient is $1+tA(0)$.
The displayed description of cycles requires this coefficient to be divisible by $t$ in $R_2$.
But $A(0)\in R_2$, and $1+tA(0)$ has value one at $t=0$.
This contradiction proves the proposition.
\end{proof}

Thus an all-arity differential-module comparison cannot retain this regular coefficient-compatible connection on its binary restriction.
It must specify a different source enlargement, a different coefficient action, or different translation relations.
The nonadjacent extension supplies the ternary source term and its exact cohomology cost.
Its compatibility with binary source replacement, ordered surjections, deconcatenation, and subsequent collisions remains a separate construction.

\begin{thebibliography}{9}
\bibitem{BDnonadjacent}
A. Beilinson and V. Drinfeld,
\emph{Chiral Algebras}, American Mathematical Society, 2004,
\S3.4.11, equation (3.4.11.1), pp.~182--183.
\end{thebibliography}
\end{document}
